\begin{longtable}{@{}p{3.0cm}p{3.5cm}p{3.5cm}p{3.5cm}@{}}
\caption{Perbandingan karakteristik pendekatan pembelajaran mesin
untuk prediksi kadar glukosa.}
\label{tab:tab2.1_perbandingan_model} \\

\toprule
\textbf{Dimensi} &
\textbf{Random Forest} &
\textbf{Gradient Boosting} &
\textbf{LSTM} \\
\midrule
\endfirsthead

\multicolumn{4}{@{}l}{\small\emph{Tabel \thetable{} (lanjutan)}} \\
\toprule
\textbf{Dimensi} &
\textbf{Random Forest} &
\textbf{Gradient Boosting} &
\textbf{LSTM} \\
\midrule
\endhead

\bottomrule
\endlastfoot

        Representasi temporal 
        &
        Memanfaatkan fitur historis yang telah direpresentasikan dalam bentuk tabular. 
        &
        Memanfaatkan fitur tabular atau ringkasan historis untuk membentuk prediktor secara bertahap. 
        &
        Mempelajari dependensi temporal secara langsung melalui urutan data. \\

        Karakteristik model 
        &
        Ensemble sejumlah pohon keputusan dengan randomisasi fitur.
        &
        Ensemble yang dibangun secara bertahap untuk memperbaiki kesalahan model sebelumnya. 
        &
        Jaringan saraf rekuren dengan mekanisme gerbang untuk mempertahankan informasi temporal. \\

        Kompleksitas
        &
        Relatif lebih sederhana pada representasi tabular.
        &
        Dapat meningkat sesuai jumlah pohon dan konfigurasi boosting.
        &
        Lebih kompleks karena melibatkan pembelajaran jaringan saraf dan sekuens.
        \\

        Kebutuhan komputasi
        &
        Umumnya lebih ringan dibandingkan pendekatan deep learning pada skenario sebanding.
        &
        Bergantung pada jumlah pohon, kedalaman, dan konfigurasi boosting.
        &
        Relatif lebih tinggi karena proses pelatihan jaringan saraf.
        \\

        Kelebihan utama
        &
        Mampu menangani hubungan nonlinear pada fitur tabular.
        &
        Mampu membentuk prediktor kuat melalui pembelajaran bertahap.
        &
        Mampu mempelajari dependensi urutan temporal secara langsung.
        \\

        Keterbatasan utama
        &
        Temporalitas tidak dipelajari secara langsung tanpa rekayasa fitur.
        &
        Bergantung pada kualitas representasi fitur temporal.
        &
        Membutuhkan data dan komputasi lebih besar serta tidak selalu menjamin generalisasi lebih baik.
        \\
        
        \noindent\tabnotestyle Sumber: Disusun penulis berdasarkan \textcite{breiman2001,friedman2001,hochreiter1997,ghimire2024}.\normalsize
\end{longtable}